\chapter{Related Work}
\label{ch:relatedwork}

Present the state of the art relevant to your research questions. This chapter should give the reader enough background to understand your contribution and demonstrate that you are aware of what has already been done.\footnote{If the related work spans multiple distinct domains, organise it into sections — one per domain. This makes the chapter easier to navigate.}

When citing, integrate author names naturally into your sentences when presenting their work, rather than appending citations as afterthoughts. For example: ``Smith et al.\ showed that \ldots~[1]''.

To find relevant publications, use academic search engines and domain-specific libraries.\footnote{%
  General: Google Scholar, Scopus, Web of Science, Semantic Scholar.
  Computer science: ACM Digital Library, IEEE Xplore, arXiv (preprints).
  Engineering \& sciences: ScienceDirect (Elsevier), SpringerLink, ASME Digital Collection, ASCE Library, PLOS.
  Natural sciences: PubMed, Web of Science.
  Open-source software: Journal of Open Source Software (JOSS), Journal of Open Research Software (JORS).
  Prefer peer-reviewed conference and journal papers over web sources. Use footnotes for web links and \texttt{\textbackslash cite} for articles, books, and papers.}

%%%%%%%%%%%%%%%%%%%%%%%%%%%%%%%%%%%%%%%%%%%%%%%%%%%%%%%%%%%%%%%%%%%%%%%%
\section{Topic Area One}
\label{sec:related-topic-one}

Summarise the key works in the first relevant area. Highlight what they contribute and where they fall short with respect to your problem.

%%%%%%%%%%%%%%%%%%%%%%%%%%%%%%%%%%%%%%%%%%%%%%%%%%%%%%%%%%%%%%%%%%%%%%%%
\section{Topic Area Two}
\label{sec:related-topic-two}

Summarise the key works in the second relevant area. End the chapter by making clear how your thesis builds on and differs from the existing literature.
